\chapter{Conclusões}
\label{cap:conclusoes}

\section{Considerações Finais}

O presente trabalho teve como motivação relatar a experiência prática de evolução e refatoração do Sistema Integrado de Processos Seletivos (SIPROS) no escopo da Fábrica de Software da Universidade Federal de Goiás. A metodologia adotada fundamentou-se na pesquisa-ação e em processos inspirados em abordagens ágeis, garantindo adaptação contínua e entregas incrementais.

A análise da vivência demonstra que o objetivo geral foi plenamente alcançado: a equipe implementou com sucesso os mecanismos de segurança necessários (Keycloak e PKCE) e consolidou práticas robustas de garantia da qualidade (QA) no \textit{front-end} e \textit{back-end}. De forma análoga, os objetivos específicos foram concretizados por meio do mapeamento de V\&V da interface, da adoção sistemática de \textit{Code Review} e da reestruturação da base de testes automatizados, erradicando abordagens provisórias herdadas e entregando um sistema arquiteturalmente mais maduro.

\section{Síntese da Experiência na Fábrica de Software}

A imersão na Fábrica de Software revelou uma disparidade substancial entre projetos estritamente acadêmicos e a atuação em um sistema de criticidade real. A responsabilidade exigida ao transacionar dados sensíveis e gerenciar a infraestrutura de uma aplicação com impacto institucional exigiu da equipe um nível de profissionalismo inerente ao mercado de trabalho, moldando positivamente a postura técnica e ética dos discentes envolvidos.

No que tange aos aprendizados, a experiência proporcionou uma evolução evidente tanto em competências técnicas (\textit{hard skills}) quanto comportamentais (\textit{soft skills}). Houve grande consolidação no domínio do \textit{framework} Angular, na orquestração de contêineres via Docker e na integração de segurança com a \textit{Google Cloud Platform} (GCP) e OAuth 2.0. Sob a ótica comportamental, o projeto instigou aprimoramentos no autogerenciamento, na comunicação interativa com os pares e na autonomia para defender e aplicar decisões arquiteturais diante de problemas complexos.

As principais dificuldades enfrentadas concentraram-se na curva inicial de aprendizado para a configuração de infraestrutura (\textit{Keycloak}) e na ausência temporária de uma "fonte da verdade" unificada para os requisitos de \textit{front-end}. 

Como lições aprendidas e sugestões para os estudantes que vivenciarão a Fábrica de Software em semestres futuros, recomenda-se fortemente:
\begin{enumerate}
    \item \textbf{Comunicação Proativa e Cultura de Code Review:} Em ambientes colaborativos, a transparência sobre qual funcionalidade está sendo desenvolvida e a disposição para revisar o código dos colegas são essenciais para evitar conflitos no repositório. O processo de revisão deve ser encarado não como um obstáculo, mas como a principal ferramenta de disseminação de conhecimento.
    \item \textbf{Gestão Ativa de Débito Técnico e Código Legado:} Não tenham receio de atuar em bases de código legadas. O primeiro instinto de um desenvolvedor iniciante pode ser reescrever todos os componentes do zero, mas a verdadeira maturidade em Engenharia de Software demonstra-se na habilidade de analisar, compreender o legado e refatorá-lo de forma incremental e segura.
    \item \textbf{Priorização da Garantia de Qualidade:} A cultura de testes de unidade nunca deve ser tratada como uma etapa secundária ou deixada para o final. Valorizem a escrita e a automação de testes desde o primeiro dia da \textit{Sprint}, garantindo que a evolução do software não se torne refém de regressões ocultas no futuro.
    \item \textbf{Definição de uma Fonte da Verdade:} Sempre certifiquem-se de alinhar os requisitos técnicos com os documentos formais (como os Casos de Uso) antes de iniciar a prototipação e o desenvolvimento no \textit{front-end}. O alinhamento precoce previne o desperdício de tempo decorrente de inúmeros retrabalhos de interface.
\end{enumerate}

\section{Trabalhos Relacionados}

O relato descrito corrobora com a literatura existente no campo de Engenharia de Software Empírica, particularmente naqueles trabalhos que documentam os desafios na adoção de metodologias inspiradas em práticas ágeis em ambientes educacionais controlados. As observações sobre a necessidade de rigor na transição de fluxos de autenticação em estágio de prototipação para implementações maduras baseadas no protocolo OAuth 2.0 com PKCE também encontram respaldo em estudos de caso contemporâneos sobre Segurança em Aplicações de Página Única (SPAs). Observa-se que a adoção de práticas sólidas de V\&V, testes automatizados e \textit{Code Review} continua sendo apontada pela literatura como a estratégia mais eficaz de mitigação de débito técnico orgânico.

\section{Trabalhos Futuros}

Apesar da evolução estrutural promovida neste ciclo, o SIPROS apresenta oportunidades contínuas para aperfeiçoamento. Diante das limitações temporais do semestre e do escopo do projeto, sugerem-se os seguintes trabalhos futuros:

\begin{itemize}
    \item \textbf{Testes de Ponta a Ponta (\textit{End-to-End}):} Iniciar a implementação de testes E2E no \textit{front-end} utilizando ferramentas como \textit{Cypress}, a fim de simular e validar o fluxo completo do usuário, complementando a suíte de testes de unidade já existente.
    \item \textbf{Expansão da Cobertura do \textit{Back-end}:} Dar seguimento à estruturação das rotinas de QA, visando atingir uma cobertura abrangente nos serviços e controladores de maior criticidade da aplicação.
    \item \textbf{Conclusão do V\&V de Interface:} Finalizar o mapeamento de Verificação e Validação para os componentes visuais e telas secundárias remanescentes, garantindo a paridade total de todos os módulos com os modelos textuais de Casos de Uso.
    \item \textbf{Otimização do Ambiente Local:} Avaliar a reestruturação dos contêineres Docker e a carga de serviços simulados localmente no repositório, com o intuito de reduzir o elevado consumo de recursos de \textit{hardware} (\textit{overhead}) nas máquinas de desenvolvimento pessoal.
\end{itemize}